% Generated draft: requires mathematical and source review.
\documentclass[11pt]{article}
\usepackage[margin=25mm]{geometry}
\usepackage{fontspec}
\setmainfont{Latin Modern Roman}
\usepackage{amsmath,amssymb,mathtools,unicode-math}
\setmathfont{Latin Modern Math}
\usepackage{enumitem,needspace,xurl,hyperref}
\setlength{\parindent}{0pt}
\setlength{\parskip}{5pt}
\setlength{\emergencystretch}{3em}
\newcommand{\N}{\mathbb{N}}
\newcommand{\Z}{\mathbb{Z}}
\newcommand{\Q}{\mathbb{Q}}
\newcommand{\R}{\mathbb{R}}
\newcommand{\C}{\mathbb{C}}
\newcommand{\F}{\mathbb{F}}
\newcommand{\A}{\mathbb{A}}
\newcommand{\PP}{\mathbb{P}}
\newcommand{\E}{\mathbb{E}}
\newcommand{\eps}{\varepsilon}
\newcommand{\dd}{\,\mathrm{d}}
\begin{document}
\section*{UnsolvedMath ID: 801071}
\textbf{Source problem code:} catalogue entry
\subsection{Definitions and mathematical statement}
Let \(X\) be a projective surface defined over \(\R\), and let \(f\colon X\to X\) be an automorphism defined over \(\R\), as in the supplied source. Write \(X(\C)\) and \(X_{\R}=X(\R)\) for its complex and real points, respectively, with their analytic topologies. Let \(f_{\C}\) denote the induced map on \(X(\C)\), and let \(f_{\R}=f_{\C}|_{X_{\R}}\). For a continuous self-map \(g\), write \(h_{\mathrm{top}}(g)\) for its topological entropy.

\textbf{Question 23.} For every real surface automorphism in the intended source class, does
\[
 h_{\mathrm{top}}(f_{\R})=h_{\mathrm{top}}(f_{\C})
 \quad\Longrightarrow\quad
 f\text{ is quasi-hyperbolic}
\]
hold?

\textbf{Unresolved qualifications.} The supplied record does not define quasi-hyperbolicity for these surface automorphisms or fully specify the intended class of surfaces. In particular, it does not explicitly state smoothness, connectedness, nonemptiness of \(X_{\R}\), or a positive-entropy hypothesis. These conditions have not been added. Consequently, the implication above records the intended question but requires the source's definitions and ambient hypotheses before it becomes a fully specified mathematical statement.

The source also writes
\[
 h_{\mathrm{top}}(f_{\R})\leq \log\lambda
\]
and attributes to reference [BS8] the assertion that, for generalized H\'enon maps, equality \(h_{\mathrm{top}}(f_{\R})=\log\lambda\) implies quasi-hyperbolicity. The supplied excerpt does not define \(\lambda\) or provide the bibliographic identity of [BS8]. This contextual assertion is preserved as a source report; it is not an independently verified theorem or an additional hypothesis of the question.
\subsection{Short English statement}
The question asks whether a real projective surface automorphism must be quasi-hyperbolic whenever its restriction to the real points has the same topological entropy as its action on the complex surface. The supplied excerpt leaves quasi-hyperbolicity and some ambient hypotheses unspecified. The local registry labels the problem open; this is a supplied status, not an independently verified assessment.
\subsection{Sources}
\begin{itemize}
\item Local UnsolvedMath registry, permanent ID \(801071\), source code \texttt{catalogue entry}, title ``Henon Maps: Open Problems --- Question 23.'' The supplied record attributes the problem to Berger and Bedford, with proposed year 2023, and labels its status open. Upstream revision: \texttt{d371ee917ac9647fbb06cd6fc529deec6f3dda82}.
\item Berger and Bedford, \emph{H\'enon Maps: A List of Open Problems} (2023), Question 23. Supplied source URL: \url{https://arxiv.org/abs/2312.03907}. The record reports access on 2026-10-16 and extraction by \texttt{arxiv\_latex}; these are supplied provenance details, not a claim of independent access or verification.
\item Reference [BS8], cited in the supplied excerpt for the generalized H\'enon-map result. No authors, title, publication details, or URL for [BS8] are supplied, so its bibliographic identity and precise theorem hypotheses remain unresolved.
\end{itemize}
\end{document}
